\documentclass[a4paper,12pt]{article}
\usepackage{color}
\usepackage{graphicx, gensymb}
\usepackage{amsfonts, cancel, amssymb}
\usepackage{amsmath, array, esvect, esint}
\usepackage{typearea}
\usepackage[a4paper,left=2cm,right=2cm,top=2cm,bottom=3cm,bindingoffset=5mm]{geometry}
\newcommand\numberthis{\addtocounter{equation}{1}\tag{\theequation}}
\newcommand{\bra}{}
\newcommand{\kom}{}
\newcommand{\brak}{}
\newcommand{\braket}{}
\newcommand{\intx}{}
\newcommand{\intr}{}
\newcommand{\kla}{}
\newcommand{\klam}{}
\newcommand{\klammer}{}
\newcommand{\oper}{}

\begin{document}

\title{04 Zustände im Wasserstoff-Atom}
\author{Joachim Schwardt}
\date{\today}
\maketitle

\section*{Aufgabe 10: Zustände bei $n=3$ im H-Atom}
\subsection*{a) }
Für $n=3$ sind im H-Atom folgende 5 Zustände möglich:\\ \\
\begin{tabular}{c|c|c}
$l$&$j_+$&$j_-$ \\ \hline
0 & 1/2 &   \\ \hline
1 & 3/2 & 1/2 \\ \hline
2 & 5/2 & 3/2
\end{tabular}
\subsection*{b) }
\section*{Aufgabe 11: Zusammengesetzte Drehimpulse}
\subsection*{a) }
Charakterisierung von $3d_{3/2}$ und $3d_{5/2}$ mit den Quantenzahlen $n,l,j$:
\begin{align*}
3d_{3/2}&:\ n=3,\ l=1\ \mathrm{oder}\ 2,\ j=\frac{3}{2}\\
3d_{5/2}&:\ n=3,\ l=2,\ j=\frac{5}{2}
\end{align*}
\subsection*{b) }
\subsection*{c) }
Es gilt $j^2=l^2+s^2+2ls\cos\alpha$. Für den maximal möglichen Wert von $l$ gilt dann:
\begin{align*}
3d_{3/2}&:\ \alpha = \arccos\kla\left( {\frac{j^2-l^2-s^2}{2ls}} \right)=\pi \\
3d_{5/2}&:\ \alpha= 0
\end{align*}
\section*{Aufgabe 12: Feinstruktur im Helium-Ion}
Nach \textsc{Dirac} lautet die Feinstruktur Korrektur:
$$E_{FS}=\frac{E_nZ^2\alpha^2}{n}\klam\left[ {\frac{2}{2j+1}-\frac{3}{4n}} \right]$$
\subsection*{a) }
Der Korrekturterm kann nicht verschwinden:
\begin{align*}
0&\overset{!}{=}\frac{3}{4n}-\frac{2}{2j+1}=\frac{12n(2j+1)-8n}{4n(2j+1)}\ \ \Rightarrow\ \ 12n(2j+1)\overset{!}{=}8n\ \ \Rightarrow\ \ j=-\frac{1}{6} 
\end{align*}
Es gilt $\frac{1}{2}\le j\le l+\frac{1}{2}$ und $0\le l\le n-1$. Daraus folgt $E_{FS}<0$:
\begin{align*}
E_{FS}=-\frac{|E_n|Z^2\alpha^2}{n}\klam\left[ {\frac{2}{2j+1}-\frac{3}{4n}} \right]\le -\frac{|E_n|Z^2\alpha^2}{n}\klam\left[ {\frac{2}{2n}-\frac{3}{4n}} \right]< 0
\end{align*}
\subsection*{b) }
Es ist $j=l\pm\frac{1}{2}$ und $l\le n-1$. Also ist die Anzahl an Energieniveaus für $n=3,\, 4$:
$$\# E_3 = \Big |\klammer\left\{ {j\,\Big|\, \frac{1}{2},\frac{3}{2},\frac{5}{2}} \right\}\Big |=3\ \ \ \ \mathrm{und}\ \ \ \ \# E_4=\Big|\klammer\left\{ {j\,\Big|\, \frac{1}{2},\frac{3}{2},\frac{5}{2},\frac{7}{2}} \right\}\Big|=4 $$

\subsection*{c) }



\end{document}